% Generated draft: requires mathematical and source review.
\documentclass[11pt]{article}
\usepackage[margin=25mm]{geometry}
\usepackage{fontspec}
\setmainfont{Latin Modern Roman}
\usepackage{amsmath,amssymb,mathtools,unicode-math}
\setmathfont{Latin Modern Math}
\usepackage{enumitem,needspace,xurl,hyperref}
\setlength{\parindent}{0pt}
\setlength{\parskip}{5pt}
\setlength{\emergencystretch}{3em}
\newcommand{\N}{\mathbb{N}}
\newcommand{\Z}{\mathbb{Z}}
\newcommand{\Q}{\mathbb{Q}}
\newcommand{\R}{\mathbb{R}}
\newcommand{\C}{\mathbb{C}}
\newcommand{\F}{\mathbb{F}}
\newcommand{\A}{\mathbb{A}}
\newcommand{\PP}{\mathbb{P}}
\newcommand{\E}{\mathbb{E}}
\newcommand{\eps}{\varepsilon}
\newcommand{\dd}{\,\mathrm{d}}
\begin{document}
\section*{UnsolvedMath ID: 30005810}
\textbf{Source problem code:} catalogue entry
\subsection{Definitions and mathematical statement}
For a graph \(G\), let \(V(G)\) denote its vertex set and let \(n=|V(G)|\) be its number of vertices. The supplied corrected statement asks whether there exists an absolute constant \(C>0\), independent of both \(G\) and \(n\), such that every grid graph with \(n\) vertices has a \(1\)-thick drawing contained in a square of side length \(C\sqrt n\).

The supplied record does not define either ``grid graph'' or ``\(1\)-thick drawing''. Accordingly, let \(\mathcal G_{\mathrm{grid}}\) provisionally denote the intended class of grid graphs, and let \(\mathcal T_1(G)\) provisionally denote the intended collection of \(1\)-thick drawings of \(G\). For a drawing \(D\), write \(\operatorname{im}(D)\) for its geometric image. With these unresolved conventions, the question for positive vertex counts has the form
\[
\exists C>0\;\forall n\in\Z_{\geq 1}\;\forall G\in\mathcal G_{\mathrm{grid}},\qquad
|V(G)|=n\ \Longrightarrow\ 
\exists D\in\mathcal T_1(G)\quad
\operatorname{im}(D)\subseteq[0,C\sqrt n]^2.
\]
Here the coordinate square represents the square in the question. No additional clearance from its boundary is specified in the supplied record.

\textbf{Unresolved definitions.} The record does not specify whether grid graphs are arbitrary finite subgraphs of a lattice graph, induced subgraphs, or rectangular grid graphs, nor does it specify the lattice dimension. It also does not specify the permitted edge geometry, intersection rules, or distance conditions defining thickness, including how incident vertices and edges are treated. The notation above therefore records the quantifiers and size bound without supplying these missing definitions; it is not a complete intrinsic formulation until those conventions are established. The source does not separately address the empty graph.

The bound \(C\sqrt n\) follows the supplied corrected statement. The original extraction has a displaced square-root symbol and extensive unrelated text; it is not a reliable independent rendering of that bound.
\subsection{Short English statement}
Is there a universal constant \(C>0\) such that every grid graph on \(n\) vertices has a \(1\)-thick drawing inside a square of side length \(C\sqrt n\)? The intended meanings of grid graph and thickness must be established before this becomes a fully specified mathematical question.
\subsection{Sources}
\begin{itemize}
\item Supplied local problem registry: permanent UnsolvedMath ID \(30005810\), title \emph{Thick Drawings of Grid Graphs}, exact source code \texttt{catalogue entry}, in the Oberwolfach Reports Open Problems collection; proposed year recorded as 2024. No proposer is identified in the dedicated attribution field.
\item Supplied report reference and literature-source link: \emph{Discrete Geometry} (2024), Oberwolfach Reports, pp.\ 137--202, DOI 10.4171/owr/2024/3, \url{https://doi.org/10.4171/owr/2024/3}. The registry labels the statement \texttt{corrected\_verified} and says it was reconstructed and checked against the official report. This is supplied verification metadata, not an independent verification performed for this draft.
\item The original extraction says the question is based on the ``Sponge problem'' and gives the incomplete reference: Larry Guth, \emph{The Width-Volume Inequality}, Geometric and Functional Analysis 17 (2007), with the page information missing. No separate URL is supplied for this reference.
\item The corrupted extraction includes the heading \emph{Open Problems in Discrete Geometry Collected by Nikita Gladkov}. Its relation to the target problem requires confirmation. Names attached to subsequently extracted Problems 2 and 3 do not establish authorship of this problem.
\item Supplied literature assessment, checked on \texttt{2026-08-22}: the problem is labeled open and described as explicitly listed as open; the assessment states that existing thick-drawing results for special graph families do not establish a universal \(C\) for every grid graph. This is supplied dated curation derived from the OWR record, not an independently verified literature assessment. The supplied background requests independent verification of time-sensitive status and cited claims.
\end{itemize}
\end{document}
